\chapter{Literature Survey}

\section{Overview of JSON and JSON Parsing}

JavaScript Object Notation (JSON) is a lightweight, text-based data format mainly used for data transfer in modern applications. It shows data using key-value pairs, arrays, and nested objects, developing it both human-readable and easy for machines to parse~\cite{ietf_rfc9535}.

JSON is mainly used in REST APIs, configuration of files, and client-server communication. Efficient parsing of JSON is impotent for the performance and reliability of software systems~\cite{sql_json_wiscorp}.

A JSON parser reads JSON text and transfer it into internal data structures like dictionaries or arrays. Parsers strictly follow the JSON syntax defined in RFC 8259 to ensure that it is correct or not and interoperability~\cite{ietf_rfc9535}.

\section{Existing JSON Parsers}

Many JSON parsers exist for many programming languages, works for both basic parsing and advanced querying features such as JSONPath:

\begin{itemize}
	\item \textbf{Built-in parsers}: Many languages include native JSON parsing, such as \texttt{JSON.parse()} in JavaScript, \texttt{json.loads()} in Python, and Java’s \texttt{JSON.parse()}~\cite{programmerall_article}.
	\item \textbf{Database parsers}: Modern databases like PostgreSQL support JSON and JSONB data types and gives built-in JSONPath query parsing and execution capabilities~\cite{postgres_jsonpath}.
	\item \textbf{Library-based parsers}: Libraries like JsonCons.Net for C\# and Jackson for Java offer rich JSON parsing and JSONPath querying support~\cite{jsoncons_net}.
	\item \textbf{Node.js JSONPath libraries}: In JavaScript server environments, libraries like \texttt{jsonpath} by Bruno Sciessere (brunerd/jsonpath) gives flexible and powerful JSONPath querying functionality~\cite{brunerd_jsonpath}. These makes complex JSONPath expressions to be evaluated efficiently in Node.js applications.
	\item \textbf{Custom and grammar-based parsers}: Tools like ANTLR4 gives grammar specifications to develop parsers capable of processing JSON and JSONPath queries~\cite{antlr4_jsonpath}.
\end{itemize}

These parsers differ in performance, feature completeness, and error-handling.

\section{Parsing Techniques}

Parsing JSON typically involves two key stages:

\begin{enumerate}
	\item \textbf{Lexical Analysis (Tokenization)}: separate input JSON text into tokens like strings, numbers, brackets, and punctuation.
	\item \textbf{Syntactic Analysis (Parsing)}: compare tokens with grammar rules to generate data structures.
\end{enumerate}

Common parsing technics:

\begin{itemize}
	\item \textbf{Recursive Descent Parsing}: it is a top-down method and parsing functions call themselves to handle nested data structures. Applicable for JSON’s with simple grammar~\cite{jsonpath_plus}.
	\item \textbf{LR Parsing}: it is a bottom-up parsing approach for parser generators like Bison or Yacc, which reduce input tokens with grammar rules to build parse trees~\cite{postgres_jsonpath}.
	\item \textbf{Parsing Expression Grammars (PEG)}: Available in some modern parsers for their expressive power and User-friendliness. ~\cite{antlr4_jsonpath}.
\end{itemize}

JSONPath parsers use these techniques to handle expressions like filters, wildcards, recursive descent selectors, and array slicing, requiring enhanced grammar and evaluation logic~\cite{ietf_rfc9535}.

\section{Challenges in JSON Parsing}

\begin{itemize}
	\item \textbf{Handling Deep Nesting}: It is possible that JSON objects and arrays is deeply nested and it is require for parsers to accurately manage recursion and memory consumption~\cite{toolsqa_jsonpath}.
	\item \textbf{Error Detection and Recovery}: Parsers need to detect malformed JSON or invalid JSONPath expressions and recover where possible~\cite{jsonpath_com}.
	\item \textbf{Performance}: Large JSON documents and complex queries need parsers optimization for speed and low memory consumption~\cite{programmerall_article}.
	\item \textbf{Supporting Advanced Query Features}: developing complex filters, logical expressions, array slicing, and recursive descent in JSONPath requires a high degree of complexity for grammar design and evaluation strategies~\cite{ietf_jsonpath_draft}.
	\item \textbf{Standards Compliance}: Parsers need to act according to formal specifications like RFC 8259 (JSON) and RFC 9535 (JSONPath) to make sure interoperability across systems~\cite{ietf_rfc9535,ietf_jsonpath_draft}.
\end{itemize}

\section{Summary}

In summary, JSON is a basic data format which need complex and affective parsing solutions. There is many parsers starting from simple language built-ins to advanced libraries and database engines supporting complex JSONPath querying~\cite{brunerd_jsonpath,jsoncons_net,postgres_jsonpath}.

Parsing methods like recursive descent and LR parsing gives the foundation for reliable JSON and JSONPath processing. However, challenges like deep nesting, error handling, performance, and supporting complex queries still remain important~\cite{toolsqa_jsonpath}.

To understand existing parsers, their techniques, and the available challenges gives deep understanding for designing efficient and robust JSON and JSONPath parsers.
